\section{From a matched local field to the raw OPE}
\label{sec:ope-grid}
The raw comparison in Input~\ref{input:ope-raw-estimates} is obtained in
three steps. At a fine lattice configuration the product is replaced by a
single local field of a sufficiently high degree. The finite tests defining
the prescription then identify its low-degree coefficients. Finally,
derivatives with respect to the cluster positions extend the estimate from
the lattice to every separated configuration. The last step uses the
filtration: taking a derivative increases the field degree by one.

We use upper rates as in Section~\ref{sec:ope-estimates}, with any sequence
$t_n\to\infty$. A cluster $X_n=(x_{i,n})_{i\in I}$ is at scale $x>0$
if $\max_i|x_{i,n}|$ has upper rate $-x$ and, for every $\delta>0$,
\[
 |x_{i,n}-x_{j,n}|>e^{-(x+\delta)t_n}\quad(i\ne j)
 \quad\hbox{eventually}.
\]
These conditions tolerate factors with zero logarithmic rate. In particular,
they include every cluster of radius $Mr_n$ and separation $\eta r_n$,
where $t_n=-\log r_n$. Spectators in this section remain a fixed positive
distance from the origin and one another, although their positions and
phases can vary with $n$.

\subsection{Intermediate scales and generic grids}
Write $a_k$ for the mesh and $s_k=\log(1/a_k)$ for its logarithmic scale.
The following elementary consequence of the mesh schedule will be useful.
Suppose $s_k$ is strictly increasing, tends to infinity, and
$s_{k+1}/s_k\to1$. Then for every $d>0$ and every $t_n\to\infty$
there are integers $j_n\to\infty$ such that $s_{j_n}/t_n\to d$.
Indeed choose the first $j_n$ with $s_{j_n}>dt_n$. Eventually $j_n>0$,
and
\[
 s_{j_n-1}\leq dt_n<s_{j_n},\qquad
 1<\frac{s_{j_n}}{dt_n}\leq\frac{s_{j_n}}{s_{j_n-1}}\longrightarrow1.
\]
The zero-rate variant follows by replacing $dt_n$ with $\sqrt{t_n}$.
The gentle-mesh condition used in the construction gives exactly the
ratio property, since $s_{k+1}-s_k=\log L_k=o(s_k)$.
The proof here uses that condition; its production for the chosen
trajectory remains part of the construction estimates.

A configuration is \emph{coordinate-generic} if, for each of the four
coordinates separately, the coordinates of its distinct points are
pairwise different. Such configurations can approximate arbitrary ones
at the mesh scale with a bound independent of their location. To see
this, enumerate $I$ by $0,\ldots,N-1$, set $K=N+1$, and for mesh $h>0$
put
\[
 y_{i,\mu}=h\left(K\left\lfloor\frac{x_{i,\mu}}{hK}\right\rfloor+i\right),
       \qquad \mu=1,\ldots,4.
\]
Different point indices have different residues modulo $K$, so each
coordinate is injective. The floor inequality and $0\leq i<K$ give
$|y_{i,\mu}-x_{i,\mu}|\leq hK$, hence
\begin{equation}\label{eq:ope-grid-round}
 \max_i|y_i-x_i|\leq2Kh.
\end{equation}
This is the empty-exterior version of the grid construction; no claim is
made that it preserves coordinate distances from fixed exterior points.
Only ordinary Euclidean separation of the full configuration is needed
below.

If $h_n=e^{-dt_n+o(t_n)}$ with $d>x$, the rounded cluster and every
configuration in the closed ball
\[
 B_n=\{Z:\|Z-X_n\|\leq\|Y_n-X_n\|\}
\]
remain at scale $x$, using the maximum of the Euclidean norms on the
configuration space. Here is the separation check. Given $\delta>0$,
choose $0<\epsilon<\delta$ with $x+\epsilon<d$ and use the original
separation at exponent $x+\epsilon$. The displacement has strictly
smaller exponential order, so eventually twice its norm is at most half
that separation. Thus the perturbed separation is at least
$\tfrac12e^{-(x+\epsilon)t_n}>e^{-(x+\delta)t_n}$. The radius bound
follows from the triangle inequality. Fixed exterior clearances also
persist, because these clusters tend to the origin.

\subsection{Reading a high-degree match}
We state the still needed quantitative local matching input precisely.
All fields and tests in it belong to the same extracted hierarchy and
the prescription already constructed in Section~\ref{sec:ope-prescription}.
For a cluster label list $p$, write $w(p)=\sum_i\deg p_i$ and
\[
 \tau(p;a,d)=d\,w(p)+4|I|(d-a).
\]
The second term is the four-dimensional volume cost of the local
matching estimate.

\begin{assumption}[Local matching estimates]\label{input:ope-match}
Fix $0<b<a<x<d$, a nonempty cluster label list $p$, a
coordinate-generic configuration sequence $X_n$ on meshes of logarithmic
rate $d$, at cluster scale $x$, and an integer $q\geq0$.
There are degree-$q$ local labels $Q_n$, and nonnegative budgets $A_n,V_n$,
with upper rates
\[
 A_n:\quad\tau(p;a,d)-(q+\tfrac1{16})(a-b),\qquad
 V_n:\quad\tau(p;a,d)-qa,
\]
with the following properties. The discrepancy between the cluster
product and $Q_n$ inserted at the origin, tested against any fixed
spectator label list at fixed positive clearance, has upper rate at
most that of $A_n$, uniformly along every choice of phase and exterior
configuration. The same is true for the mixed state and the finite tests
of the prescription. The labels obey
\begin{equation}\label{eq:ope-match-norm}
 \|Q_n\|\leq J_n V_n,\qquad J_n\text{ has upper rate }qa.
\end{equation}
For each $1\leq m\leq q$ there are $z_{m,n}$ in the label space of
degree $q$ whose physical classes lie in the included space $E_{m-1}$,
and
\begin{equation}\label{eq:ope-match-distance}
 \|Q_n-z_{m,n}\|\leq F_{m,n}V_n,
 \qquad F_{m,n}\text{ has upper rate }(q-m)a.
\end{equation}
All norms here are fixed finite-dimensional label norms. The original
labels $Q_n$, before passing to their physical classes, satisfy these
bounds. The estimates hold for every fixed cluster label list, including
lists formed by differentiating its fields.
\end{assumption}
In the source the same $Q_n$ arises as a finite-dimensional sublimit of
cutoff matches, retaining their local discrepancy and quota bounds.
The individual channel estimate below is essential: a bound on the
whole degree-$q$ label would lose the powers that distinguish low and
high channels. Proposition~\ref{prop:ope-match-assembly} derives
Input~\ref{input:ope-match} from the cutoff construction of
Section~\ref{sec:ope-cutoff}, the phase transfer of
Section~\ref{sec:ope-phase-prices} and the filtration estimates of
Section~\ref{sec:ope-filtration}, with one label $Q_n$. That derivation
is conditional on the analytic inputs listed in the proposition; the
Input is not yet derived unconditionally from the finite Wilson law.

Let $f_{\alpha,n}$ denote the raw coefficient of the product and let
$v_{\alpha,n}=b_\alpha([Q_n])$ be the coefficient of its matched field,
where $b_\alpha$ is the fixed functional of the prescription and
$m=|\alpha|\leq q$. Then
\begin{equation}\label{eq:ope-match-grade}
 v_{\alpha,n}\text{ has upper rate }\tau(p;a,d)-ma.
\end{equation}
For $m=0$, apply the fixed linear functional to
\eqref{eq:ope-match-norm}. For $m>0$, coherence of the prescription
implies $b_\alpha([z_{m,n}])=0$: this detector kills every field of degree
less than $m$. Hence
\[
 |v_{\alpha,n}|=|b_\alpha([Q_n-z_{m,n}])|
       \leq \|b_\alpha\circ\pi_q\|\,\|Q_n-z_{m,n}\|,
\]
where $\pi_q$ is the quotient map. Combining
\eqref{eq:ope-match-distance} with the rate of $V_n$ proves
\eqref{eq:ope-match-grade}.

Each raw coefficient is a finite linear combination of exterior tests.
Applying the discrepancy part of Input~\ref{input:ope-match} to those
tests therefore gives
\begin{equation}\label{eq:ope-match-detector}
 f_{\alpha,n}-v_{\alpha,n}\text{ has upper rate }
                  \tau(p;a,d)-(q+\tfrac1{16})(a-b).
\end{equation}
For a specified output degree $m$, choose $q\geq m$ so large that
$(q+\tfrac1{16})(a-b)\geq(m+1)a$. Then
\eqref{eq:ope-match-grade} and \eqref{eq:ope-match-detector} show that
$f_{\alpha,n}$ has upper rate $\tau(p;a,d)-ma$.

The analogous remainder argument retains the entire product correlation.
Let $\ell_n$ be the functional obtained by inserting a field at the origin
beside the chosen spectators. Assume, as furnished by the separated
correlation estimate, that $\ell_n(u)$ has upper rate zero for every fixed
field $u$ and every phase sequence. To truncate at $D$, choose $q\geq D$
with $(q+\tfrac1{16})(a-b)\geq(D+1)a$. Exact reconstruction of $[Q_n]$
by its prescription coefficients gives
\begin{align}
 S_{\chi_n}(p,z;X_n,Y_n)-\ell_n(H_D(X_n,p;0))
 &=\bigl[S_{\chi_n}(p,z;X_n,Y_n)-\ell_n([Q_n])\bigr]\nonumber\\
 &\quad+\sum_{|\alpha|\leq D}(v_{\alpha,n}-f_{\alpha,n})
                            \ell_n(u_\alpha)\nonumber\\
 &\quad+\sum_{D<|\alpha|\leq q}v_{\alpha,n}\ell_n(u_\alpha).
 \label{eq:ope-match-remainder}
\end{align}
The first term and each low-channel difference have upper rate at most
$\tau(p;a,d)-(D+1)a$. Each high channel has that rate as well, by
\eqref{eq:ope-match-grade}, because $a>0$ and $|\alpha|\geq D+1$.
The sums are finite. This proves the raw grid remainder estimate with
upper rate $\tau(p;a,d)-(D+1)a$. The matching label may depend on $q$;
the raw coefficients and the truncation on the left do not.

\subsection{Interpolation by descending through the grading}
We give the interpolation argument explicitly. Suppose functions
$h_n(p,X)$ are indexed by finite field lists of the same length, with
integer total degree $w(p)$. Assume that for every scale-$x$ sequence:
\begin{enumerate}
\item $h_n(p,X_n)$ has a rough upper rate $xw(p)+c$;
\item on every coordinate-generic grid sequence of logarithmic mesh rate
$d>x$, it has upper rate $dw(p)+e$;
\item it is eventually differentiable along each such sequence, with identity
\[\partial_{i,\mu}h_n(p,X_n)=h_n(p',X_n).\]
Here $p'$ is obtained by differentiating the corresponding field, so
$w(p')=w(p)+1$.
\end{enumerate}
The constants $c,e$ are independent of $p$; they may depend on fixed
truncation and exterior data. Then at every scale-$x$ sequence the
improved upper rate is $dw(p)+e$.

Choose an integer $N\geq0$ with $(d-x)N\geq c-e$. We prove by induction
on $k\geq0$ that the improved rate holds whenever $N\leq w(p)+k$.
For $k=0$ the rough rate is already no larger than the improved one.
For the induction step round $X_n$ to $Y_n$ by
\eqref{eq:ope-grid-round}. Every choice $Z_n\in B_n$ stays at scale $x$.
The induction hypothesis applied to the differentiated lists shows that
each coordinate derivative at $Z_n$ has upper rate
$d(w(p)+1)+e$. In the fixed finite-dimensional configuration space,
\[
 \|Dh_n(p,Z_n)\|\leq\sum_{i,\mu}
                       |\partial_{i,\mu}h_n(p,Z_n)|,
\]
so the derivative norm has that rate too.

These bounds and differentiability hold eventually uniformly on $B_n$.
For any specified positive slack, failure would give infinitely many
indices with a bad point; choosing a bad point there and the center
elsewhere produces a sequence $Z_n\in B_n$ contradicting the corresponding
eventual assertions. The mean-value inequality on the convex ball thus
gives
\[
 |h_n(p,Y_n)-h_n(p,X_n)|
 \leq\sup_{Z\in B_n}\|Dh_n(p,Z)\|\,\|Y_n-X_n\|.
\]
The displacement has upper rate $-d$, so the difference has upper rate
$dw(p)+e$. The grid value has the same rate, which proves the induction
step. Taking $k=N$ covers every field list. Only finitely many derivative
levels are used for any one application, although the assumptions must be
available at every finite degree.

For clarity, the mean-value inequality here needs no continuity of the
derivative. On each segment in $B_n$, differentiability implies
continuity, and the scalar mean-value theorem applied to real projections
of the complex-valued function bounds its increment by the derivative
bound times the segment length. The finite-dimensional operator estimate
above follows by expanding a vector in its coordinate basis and applying
the triangle inequality.

\subsection{Removing the intermediate-scale losses}
The differentiability identity (3) needed for this interpolation is the
slot derivative identity \eqref{ex-raw-slot-identity}: one grade-raising
choice of raw derivative representative $\delta_jv$ for every raw label,
whose slot identity holds at distinct points with arbitrary finite lists
of exterior insertions, and, because $h_n$ is a phase discrepancy, in
every phase table used here, with the chosen representatives having the
indicated values in that phase. Section~\ref{ex-slab-regularity} assumes
the same identity and does not produce it; it is a named input of the
present section. The identities persist under
finite linear combinations of detector tests and under subtraction of
the finite raw truncation. At a separated configuration its neighborhood
remains separated, so the totalized values away from that neighborhood
play no role in differentiation.

Suppose the rough separated correlation estimates give upper rate $xw(p)$
for the cluster with fixed-clearance spectators, uniformly in the phase.
They give the same rate for each raw coefficient, since its detector list
is fixed and finite. A finite truncation also has rough rate $xw(p)$:
its coefficient has that rate and each fixed single-field spectator
insertion has rate zero. Therefore both the raw coefficient and the raw
discrepancy satisfy the rough interpolation bound with $c=0$.

For a coefficient of degree $m$, the grid bound just proved is
\[
 dw(p)+4|I|(d-a)-ma.
\]
Use interpolation with $e=4|I|(d-a)-ma$ to obtain this bound at all
scale-$x$ clusters. For a remainder through degree $D$, use instead
$e=4|I|(d-a)-(D+1)a$. These offsets do not change when a cluster field is
differentiated, as required by the induction.

Now set $a=x-\epsilon$, $d=x+\epsilon$, and $b=a/2$, with
$0<\epsilon<x$. The mesh schedule supplies all three intermediate levels.
The resulting coefficient upper rate is
\[
 (w(p)-m)x+\epsilon\bigl(w(p)+8|I|+m\bigr),
\]
and the remainder upper rate is
\[
 (w(p)-D-1)x+\epsilon\bigl(w(p)+8|I|+D+1\bigr).
\]
Given any desired positive slack above the limiting exponent, choose
$\epsilon$ small enough that the displayed excess is less than half the
slack, then use the remaining half in the upper-rate bound. This proves
rates $(w(p)-m)x$ and $(w(p)-D-1)x$, respectively. No uniformity of the
intermediate constructions as $\epsilon\to0$ is needed.

Translating the cluster center to the origin preserves all distances and
the complete correlation. The raw coefficients use relative positions,
so the same estimates hold at arbitrary centers. If the cluster is empty,
its product is the identity. Exact reconstruction of that degree-zero
field makes every raw truncation exact; positive-degree detectors
annihilate it, and degree-zero coefficients are constant. Thus this case
also satisfies the stated rates. Taking $x=1$ proves
Input~\ref{input:ope-raw-estimates}, conditional on
Input~\ref{input:ope-match}, the rough separated correlation estimates,
and the slot derivative identity \eqref{ex-raw-slot-identity} with the
phase compatibility described above. The remaining task is
to produce those matching and correlation estimates for the actual
Wilson trajectory.
